\section{Completion of the cylinder estimate}\label{act:section}

We now finish the proof of Theorem~\ref{thm:cylinder}.
The canonical excess $E$ is nonnegative and vanishes at the equality
points specified in Proposition~\ref{names:canonical-entropy}.
Proposition~\ref{int:actions} supplies directions in which it decreases
when followed backwards in time. The final contradiction comes from a
path starting at an equality boundary along which this decrease remains
valid. The remaining difficulties are the unit-cost quadratic boundary
and intervals on which the angular trace vanishes.

All entropy limits in this section use the ordered limits of the canonical
construction.  In particular, a finite
list of macroscopic queries is fixed before the configuration, mesh, and
scale limits; a tangent gap is subsequently allowed to tend to zero.

\subsection{A line estimate at unit quadratic cost}

At unit quadratic cost the needed fixed-width time estimate is $2+3d$.
On an equality boundary it contradicts the canonical entropy upper bound;
in the strict-curvature case it gives the missing lower bound for
$\partial_aE$. The next lemma isolates this estimate when the orientation
marks lie close to a line determined by the side data. Its hypotheses
record time uncertainty, angular uncertainty along a fiber, and a mismatch
between the target line and the line obtained by resampling with the same
fiber label. Its conclusion concerns the entropy of one observation output.

\begin{lemma}[Fixed-width line estimate]\label{act:line-test}
Let $\delta\to0$.  Let $C$ be a common root name and let side data
$\mathscr Z$ determine $C$, a quadratic phase
$F=(X,B,P,Y,N,D)$, and coefficients $(\beta_F,\kappa_F)$.
The normalized phase, coefficients, and marks have size
$\delta^{-o(1)}$.  Conditional on $\mathscr Z$, draw two independent
observations with the same conditional law, whose marks satisfy
\begin{equation}
 r_i=\beta_F+\kappa_F\theta_i+O(\delta^{10-o(1)}).
 \tag{s-line}\label{eq:s-line}
\end{equation}
Suppose the following hypotheses hold, with arbitrary fixed strict
resolution slack and with exceptional probabilities tending to zero.
\begin{enumerate}[label=\textup{(\roman*)}]
\item The conditional score of $\theta_{i,p+t}$ given
$\mathscr Z,\theta_{i,p}$ is at least $dt$, for the finitely many
queries used below through depth $3$, including queries starting at $0$,
where $d>0$.
\item Each observation carries a fiber label $S_i$ such that $(C,S_i)$
determines its marks and the quadratic entries of $Q_i$ in the fiber
formulas to single-valued error $\delta^{10-o(1)}$.
\item Each observation has an output $O_i$ with these recovery properties.
The pair $(C,O_i)$ reads $\theta_i$ to depth $1$, and together with
$X_{p+t}$ reads $F_{p+t}$ whenever $0<p<p+t<1$, up to lists of
vanishing exponent.  If $1<p<p+t<2$ and $t\ll\eps,p$, it also predicts
$(\theta_i,r_i)_\eps$ and, using these coefficients and $F_p$, reads
\[
 K+\theta_{i,\eps}J-r_{i,\eps}H
\]
to depth $p+t$.  Here $H,J,K$ are the contact forms, evaluated in a
fixed positive-depth prefix frame and translated at $F_p$.
\item For each $i$ individually,
\[
 I_\delta(X_6;S_i\mid C)=o(1),\qquad
 i_\delta(X_p\mid C)=M(p)+o(1)
\]
in probability at the tested prefixes, where $M$ is deterministic,
$M(0)=0$, and $M(p)\ge g_*p$ for some $g_*>0$.
\item Draw $R$ independently from the posterior law of the side and
observation given $(C,S_i)$.  Its line satisfies
\begin{equation}
 \abs{\kappa_R-\kappa_F}\ge\delta^\chi
 \quad\text{with probability tending to one, for every fixed }\chi>0.
 \tag{mis}\label{eq:mis}
\end{equation}
\end{enumerate}
Then
\begin{equation}
 H_\delta(O_i\mid C)
 \ge d+g_*+2\min(1,2d)+d+\min(1,d+g_*)-o(1).
 \tag{line-time}\label{eq:line-time}
\end{equation}
In particular, $g_*\ge1-2d$ implies the lower bound $2+3d-o(1)$.
The finite accuracy constants $6$ and $10$ may be increased.
\end{lemma}

\begin{proof}
The angular capacities make $M$ Lipschitz.  Put $m(p)=M'(p)$ at its
Lebesgue points.  Fix the tested observation $i$, its posterior draw $R$,
and the other target observation $j$.  Write $A=X_F-X_R$ and, for a
small fixed $\eps>0$, condition on
\[
 \Omega=(\beta_F,\kappa_F,\theta_i,r_i,\theta_j,r_j,A)_\eps.
\]
Its normalized entropy given $C$ is $O(\eps)$.  In the contact frame
of the current prefix define
\begin{align*}
 L&=dP-2\beta_{F,\eps}\,dX+2\kappa_{F,\eps}\,dB,\\
 J'&=J-\kappa_{F,\eps}H,\qquad
 K'=K-\beta_{F,\eps}H,\\
 G&=(L,J',K'),\qquad
 x=K'+\theta_{i,\eps}J',\qquad
 y=J'-\tfrac12 A_\eps L.
\end{align*}
Unless specified otherwise, rates are conditional on $(C,\Omega,F_p)$;
an additional held reading is held to child depth $p+t$.
First restrict $p$ to compact subintervals of $(0,1)$ and $(1,2)$.
After $\eps$ and a positive contact-prefix depth are fixed, take
$t$ smaller than these depths and than $p$.

We first account for the cost of $\Omega$.  For each of the nested
refinement families
\begin{equation}\label{act:budget-families}
\begin{array}{ll}
 X\text{ beyond }p &\mid C,S_i,X_p\quad\text{or }C,S_j,X_p,\\
 \widetilde\theta\text{ beyond }p
   &\mid C,\mathscr Z_R,X_{F,6},\widetilde\theta_p,\\
 F\text{ beyond }p &\mid C,F_p,\Omega_0,
\end{array}
\end{equation}
where $\Omega_0$ is any fixed subtuple needed below, the integrated
expected information loss on adding $\Omega$ is at most
$H_\delta(\Omega\mid C)=O(\eps)$.  This is the conditional chain
rule applied along the refinement family.  The cumulative losses have
nonnegative Lipschitz limits, so their derivatives have the same
integral bound.

Before this addition the first family has rate $m(p)$ by stationarity.
For the second, the posterior draw has the original side-and-observation
marginal.  Conditional independence given $(C,S_i)$ and the first
stationarity hypothesis give
\[
 I_\delta(X_{F,6};\mathscr Z_R,\text{observation at }R\mid C)=o(1).
\]
Its time innovations therefore retain the lower rate $d$, by
\eqref{eq:LP}.  The same information comparison and $M(z)\ge g_*z$
show that $\abs A\ge\delta^\chi$ with probability tending to one
for each fixed $\chi>0$: the contrary event puts $X_F$ in a bounded
list of depth-$\chi$ bins predicted by $X_R$.

These budgets justify all subsequent rate errors as follows.  Fix a
strict proposed deficit and a positive lower bound for its probability.
Let $\eps\downarrow0$, discard layer points where one of the finitely
many loss densities is too large, and then choose typical remaining
points for the pathwise and mixed-prefix differentiations.  At each
such point let $t\downarrow0$.  The expected loss divided by $t$
tends to zero along this procedure; the likelihood inequalities
\eqref{eq:LP} turn it into a vanishing score error on all but a
vanishing probability.  The same argument gives a vanishing
product-likelihood cost, measured relative to $t\log(1/\delta)$,
for each planar projection test.  Conditioning additionally on a
child reading which does not use the removed subtuple costs no more:
the chain rule bounds the remaining information by that of the full
source, with the negligible list entropy of that reading added.
All rates have bounded capacities, so discarded probabilities and
layer sets contribute vanishing errors to the integrals.  This
procedure requires no uniform differentiability in $\eps$.

The output chain rule now gives
\begin{equation}
 H_\delta(O_i\mid C)
 \ge d+\int_0^1\bigl(\alpha_0(p)-m(p)\bigr)\,dp
       +\int_1^2\E|x|(p)\,dp-o(1),
 \tag{L-chain}\label{eq:L-chain}
\end{equation}
where $\alpha_0$ is the expected isotropic phase rate before adding
$\Omega$.  Indeed, conditional on fine phase, the time contribution
is at least $d$ after conditioning further on $\mathscr Z$.
Below depth $1$, adding the angular bits recovers the phase bits and
costs at most their rate $m$.  Between depths $1$ and $2$, first add
the coarse mark coefficients predicted by the output and then add
$\Omega$.  Translate the displayed reading at $F_p$ and use
\eqref{eq:s-line} to obtain $x$.  Tiling compact interior intervals
and taking limits proves \eqref{eq:L-chain}; bounded capacities allow
the omitted endpoints to tend to $0,1,2$.

Both angular fiber modes $\ell_i,\ell_j$ annihilate $G$ to the
required accuracy.  They give
\begin{equation}
 |F|\ge |G|+2m(p)-o(1).
 \tag{L-mode}\label{eq:L-mode}
\end{equation}
For completeness, add $S_i,X_p$ for the first mode, or $S_j,X_p$
for the second.  The fiber equations predict the parent and all
annihilating readings; replacing line coefficients by their
depth-$\eps$ representatives introduces error
$O(\delta^{p+\eps-o(1)})$.  First use
$dB+\theta_{i,\eps}dX$, which sees $\ell_j$, and then $dX$,
which sees $\ell_i$.  The two times are separated by every fixed
power threshold, by conditional independent sampling and their
time-score bound.  Each inversion thus costs an arbitrarily small
power, and the first budgets in \eqref{act:budget-families} supply
the two residual rates $m(p)$.

Next vary $\widetilde\theta$ at $R$, holding $\mathscr Z_R$ and
$X_{F,6}$ and starting inside a depth-$p$ time parent.  The shared
fiber label identifies the target mark with
$(\widetilde\theta,\beta_R+\kappa_R\widetilde\theta)$ to the
fine prescribed accuracy.  Equations \eqref{eq:leg} and
\eqref{eq:trans} give the response in $G$ as
\[
 (\kappa_R-\kappa_F)(2A,A^2,-\theta_i A^2)\,d\widetilde\theta,
\]
up to the coarse-coefficient error.  Here
$\beta_R-\beta_F=-\theta_i(\kappa_R-\kappa_F)$ to the same
accuracy.  This response annihilates $x,y$, whereas its $J'$ and
$L$ components are separated from zero by arbitrary power slack,
using \eqref{eq:mis} and the separation of $A$.  The parent is
predicted by the fiber formulas.  The remaining Taylor error is
$O(\delta^{2p-o(1)})$, which is finer than $\delta^{p+t}$.
The contact forms are frozen; no derivative of the line field is
taken.  Consequently the time innovation, with the second budget
in \eqref{act:budget-families}, gives
\begin{equation}
 |J'\mid x|\ge d-o(1),\qquad |L\mid x,y|\ge d-o(1).
 \tag{L-res}\label{eq:L-res}
\end{equation}

Apply the conditional projection rule of
Proposition~\ref{names:projection-rule} twice. For the source $(J',K')$,
initially retain only the line coefficients from
$\Omega$, and vary the time slope $\theta_{i,\eps}$.  Conditional
on $\mathscr Z$, the trimmed direction sublaw has exponent $d$;
integrating this unnormalized bound supplies the required direction
bound at the source parent.  The last budgets in
\eqref{act:budget-families} justify reinstating the remaining records
for the score and list tests.  If $S$ is the joint source rate,
\eqref{eq:L-res} gives $S\ge|x|+d-o(1)$, and the projection bound
$|x|\ge\min(1,S,(S+d)/2)-o(1)$ therefore gives
$|x|\ge\min(1,2d)-o(1)$.

For the second application condition on $x$ to child accuracy, retain
the target entries of $\Omega$, and vary $A_\eps$ in the projection
$y$ of $(L,J')$.  Here is the exact information comparison for this
conditioning.  Write $T$ for the target subtuple of $\Omega$,
$B_*=A_\eps$, and $D_*=(C,F_p,T)$.  Let $W_*$ be the child bin
of the frozen reading $x$, and let $Z_*$ be the planar child
source $(L,J')$.  Both are determined to negligible lists by a
sufficiently fine relative child phase name $V_*$ and $D_*$,
without using $B_*$.  The name $V_*$ may be taken at depth
$p+(1+\xi)t$ with arbitrarily small fixed $\xi>0$; the extra
capacity is $O(\xi t)$.  The conditional chain rule gives
\begin{align}
 I_\delta(Z_*;B_*\mid D_*,W_*)
 &\le I_\delta(V_*;B_*\mid D_*)
       +H_\delta(Z_*,W_*\mid D_*,V_*),
       \label{act:child-cmi}\\
 I_\delta(V_*;B_*\mid C,F_p,T)
 &= I_\delta(V_*;B_*\mid C,T)
       -I_\delta(F_p;B_*\mid C,T)+o(1).
       \label{act:child-budget}
\end{align}
For the first inequality, adjoin $V_*$ to the source, and use
$I(V_*;B_*\mid D_*,W_*)
\le I(V_*;B_*\mid D_*)+H(W_*\mid D_*,V_*)$.
The second equality uses nested phase bins, with the negligible
list error if their grids are only list-compatible.  The
cumulative mutual information on the right is nondecreasing
and bounded by $H_\delta(B_*\mid C,T)=O(\eps)$.
Consequently its loss density is one of the integrated budgets
already treated above.  Divide \eqref{act:child-cmi} by $t$,
use that budget at the retained layer points, and then let the
strict resolution slack $\xi$ tend to zero.  This proves the
needed vanishing source-record information after the child
conditioning.  In particular, $W_*$ uses the frozen target
coefficients and no $A$; coefficient replacement is by a known
translation at $F_p$ and an error
$O(\delta^{p+\eps-o(1)})$, not by an $A$-dependent reading.

For the direction bound, let the full target data include its
side and its two observations.  The names $D_*,W_*$ are functions
of these data.  Conditional on the full target data, $R$ still has
the posterior kernel given $(C,S_i)$.  Stationarity and
$M(z)\ge g_*z$ trim its angular law to unnormalized ball bounds
of exponent $g_*$.  More explicitly, for each required fixed depth
and strict tolerance retain only posterior angular cells whose
conditional mass given $(C,S_i)$ is at most the corresponding
$\delta^{g_*z-\text{tolerance}}$.  Their union has probability
tending to one by the stationary scores, and their unnormalized
mass in any constant-multiple depth-$z$ ball has the same bound
up to a fixed neighbor factor, for every conditioning value.
Translation and reflection give the same bound for $A$ for every
full-target value.  Integrating these restricted kernels over
the full target conditional on $(D_*,W_*)$ preserves the bound.
No normalization of the retained sublaw is used in this step.
The projection uses layers $t\ll\eps$, so rounding
the direction is harmless.  The last budget, now with mixed-prefix
linearity for the source conditioned on $x$, again permits the full
record in the test.  Its residual beyond $y$ is at least $d-o(1)$.
Thus
\begin{equation}\label{act:line-boosts}
 |x|\ge\min(1,2d)-o(1),\qquad
 |G|-|x|\ge d+\min(1,d+g_*)-o(1).
\end{equation}
Each assertion means that a fixed strict deficit of positive
probability contradicts the corresponding finite projection test
with the loss order specified above.

Finally $\alpha_0\ge\E|F|$, and
$\int_0^1m(p)\,dp=M(1)\ge g_*$.  Substitute
\eqref{eq:L-mode} and \eqref{act:line-boosts} into
\eqref{eq:L-chain}, and then send all the budget and endpoint errors
to zero.  This gives \eqref{eq:line-time}.  To check the last claim,
if $d\le1/2$, use $g_*\ge1-2d$ and
$\min(1,d+g_*)\ge1-d$; if $d\ge1/2$, use $g_*>0$ and
$\min(1,d+g_*)\ge d$.  In either case the stated lower bound is
at least $2+3d$.
\end{proof}

\subsection{Boundary applications of the line estimate}

The line estimate has two uses. At the stationary equality boundary it
forces the radial spreading needed by the interior estimates. In the
strict-curvature case it gives a fixed-width time bound directly. Both
applications must verify the posterior-line mismatch in \eqref{eq:mis}.

\begin{lemma}[The stationary equality boundary]\label{act:stationary-boundary}
In the classical equality configuration $q=c_2=1$, $J=0$, the
exact-side quadratic replica construction has the radial spreading
property at almost every strong time $a$, where $p_*(a)>b_1(a)$.
This holds on interior patches $v>0$ with a common static depth
$b(a)+v<x_*<p_*(a)$.
\end{lemma}

\begin{proof}
Here $u=0$ and $b_1=b_2=a-1$.  Use the unclipped quadratic template
and request the canonical profiles also on the fixed boundary $v=0$,
where $E(a,0)=0$.  Countably many depth queries suffice.  Choose a
typical strong $a$ at which the prediction slack persists immediately
to the left and \eqref{eq:ng} holds on that boundary for all fixed
accuracy demands.  The reached label $\lambda_a$ knows the boundary
name.  The strengthened version of \eqref{eq:mem} with $L=1$, ending
at $a$, together with \eqref{eq:Sp}, therefore gives angular prefix
scores with lower slope $k_0+2\Gamma$, to arbitrary strict slack.
The comparisons can use nearby actual gates with their earlier
errors already tending to zero.  If this slope exceeds the angular
capacity there is already a contradiction; otherwise it supplies a
positive $g_*$ with $g_*\ge1-2d$.

Fix the exact side $Z=(y_0,[P]_{x_*},t_a,R_a)$.  It is the same
side, with the same marks, at $v=0$ and on the tested interior
patch.  Suppose the required uniform power avoidance fails.
Then, along $z=\lambda h\downarrow0$, a line predicted from $Z$
has tube mass $s^{o(z)}$ at normalized width $s^z$.  This formulation
allows $\lambda$ subsequently to decrease, as required in the
high-set radial estimate \eqref{eq:rad}.  A vertical line, or a
slope of size at least $s^{-\chi z}$, restricts $\theta$ on the
bounded mark support to an interval of positive power width; the
conditional time-list bound excludes such mass.  Hence, with
subpower normalized coefficients, the remaining line event is
\[
 \abs{\beta_1(T_1)-\beta_*(T_1)}
   \le s^{b(a)+(1-o(1))z},\qquad
 \beta_*(T)=R_a+s^{b(a)}
       \bigl(\beta_F+\kappa_F(T-t_a)s^{-a}\bigr).
\]
Include successful lists of size $s^{-o(z)}$ for $[P]_{x_*}$ given
$(y_0,t_a)$, using strong prediction and the deterministic profiles.
Put $h'=z/400$, $\delta=s^{h'}$, and restrict the single final law
to this event.  All fixed strict time-score failures were power-small
in these units before restriction, by the causal time-list count;
likelihood comparison therefore preserves the time-score hypotheses.

Use $\mathscr Z=Z$, $C=C_*(a,0)$, and the representative phase at
$v=0$.  The predicted line satisfies \eqref{eq:s-line}.  Let $S_i$
contain depth-$12$ marks and representative quadratic fiber entries.
They are read to negligible lists from
$(C,\mathcal C(a+40h',0))$.  Actual and representative roots are
mutually predicted to negligible lists: guess neighboring derivative
bins if necessary and transport normal centers with the base-cell
accuracies; the true derivative varies by only $O(s^{x_*})$ on
this side.  The deterministic scores persist under the
$s^{o(h')}$ restriction.  Transferring those scores first, and then
using the entropy chain rule, transfers \eqref{eq:ng} to
$I_\delta(X_6;S_i\mid C)=o(1)$ with the preceding angular lower
profile.  This does not transfer an arbitrary small mutual
information bound directly through a rare-event restriction.
Further subsequences supply $M$.  The output
$O_i=\mathcal C(a+h',0)$ has the recovery properties of
Lemma~\ref{act:line-test}, by the displayed-shear estimates underlying
\eqref{eq:time-rate}.  Boundary equality in \eqref{eq:canon} gives
\[
 H_\delta(O_i\mid C)\le2+3d-2\Gamma-\eta+o(1).
\]
All profile errors are taken to be $o(h')$ before $z$ decreases.

It remains to verify \eqref{eq:mis}.  Its failure at a fixed small
$\chi>0$ gives, by independent posterior sampling and selection of
a most popular interval, a prediction of a $\delta^\chi$ interval
for $\kappa_F$ from $(C,S_i)$ with nonvanishing success.  Mark an
end gate near $j=a+60h'$ using the Borel inner completion mark of
Lemma~\ref{cyl:borel-paths}, inside the current end-state support,
and apply group extraction on $[a,j]$.
The successful completions have mass $s^{o(h')}$.  The label
$\lambda_j$ knows the canonical names through $a+40h'$ at width
$0=u(j)$, and hence $C,S_i$, to negligible lists.  Starting rows
know the lists for $Z$ along these paths, including the successful
derivative lists.

Choose $0<\omega<\chi$ and bin the value and slope of the predicted
trajectory at $t_a$ at widths
\[
 s^{b(a)+\omega h'},\qquad s^{b(a)-a+\omega h'}.
\]
Include the queried $t_a$ in the list.  The slope bin is predicted
at the end from $(C,S_i)$.  Once it is chosen, the value bin is read
to small lists from $\beta_j(T_j)$: comparison with $\beta_1(T_1)$
and transport from $T_j$ to $T_1$ cost $s^{b(j)-o(h')}$, and the
line-tube error is still smaller.  Extrapolation of the binned slope
back to $t_a$ costs only its chosen value width.  Disjointize the
finite possible-bin relations as in Lemma~\ref{cyl:borel-paths}.
Each retained marked end row has a witnessing completion in its
chosen bin; path concatenation puts all its predecessors in the
corresponding Borel outer starting list. This predecessor projection
uses the unrestricted truncated \([a,j]\) geometric correspondence,
not only successful completions. Its Borel inner partner carries
the same full analyzed pre-synthesis input state in the original bin
partition.
After the common subtraction, the end curvature value is bounded
by $s^{b(a)+\omega h'-o(h')}$, and its duration times the residual
slope is at most $s^{b(j)-o(h')}$.  As $b(j)-b(a)=60h'$, the
curvature cost on this segment is strictly less than $1$.
Group extraction contradicts its minimality.  Fixed list slacks
are chosen first to dominate mesh and closure errors and then sent
to zero in $h'$ units.  Only the curvature cost is improved.

Thus \eqref{eq:mis} holds.  Lemma~\ref{act:line-test} contradicts
the displayed boundary upper bound.  This proves uniform line
avoidance and hence the stated radial spreading.  The restriction
was made on the deterministic-profile single law before the new
line experiment, so no stationarity claim on a subpower part of
an already constructed many-replica law is needed.
\end{proof}

\begin{proposition}[Strict affine boundary]\label{act:strict-boundary}
Suppose $c_2=1$ and $q<1$, with no lower bound on $q$.  Use the
artificial quadratic template ending at $D$, and set $v_0=u(D)$;
in the hybrid case take $D=1$.  If the boundary prefix memory
supplies $g_*>0$ with $g_*\ge1-2d$, then at almost every $a<D$
\[
 \partial_a\mathcal H(a,v_0)\ge2+3d,
 \qquad \partial_a E(a,v_0)\ge2\Gamma.
\]
The same conclusion holds for the clipped template on this boundary.
\end{proposition}

\begin{proof}
The boundary is an open quadratic branch strictly before $D$:
\[
 l(a,v_0)=b_1(D)+a-D<\min_{[a,D]}b_1,
 \qquad w(a,v_0)=n(D)+a-D.
\]
In hybrid applications this also supplies strict normal slack.  The
derivative to be proved is at fixed $v_0$, so interior stationarity
in width is unnecessary.

We first specify the forward profiles needed in the hybrid case.
At each fixed rational $0<j<1$, preflatten the actual read rows
$\mathbf P^j$ in chronological order, before the final deterministic
profiles.  Make finitely many requests and then diagonalize.  On
these rows use a version $\mathcal C^{[j]}(a,v)$, $a<j$, with the
same exponents and endpoint constant as the artificial template,
but with the $j$-label and its coordinates as the actual data.
Namely, set $\zeta=b_2(D)-D$, bin $\bar k_j=[K_j]_\zeta$,
subtract that trajectory, and use the residual
$\beta'_j(T_j),P_j$ for the curvature and derivative inputs.
Here the prime denotes residual coordinates.  The queries include
the roots at $v_0=0$ and their angular refinements.  In the classical
case any fixed $a<j<D$ is allowed: evaluate these names via the
encountered label $\lambda_j$ on the final read law, with no change
to an earlier gate.

On connections from $j$ to $D$, the names $\mathcal C^{[j]}$ and
$\mathcal C$ have bidirectional lists of vanishing exponent,
also at nearby times below $j$.  To verify this, their global slope
bins differ by at most $s^{\zeta-o(1)}$, by \eqref{eq:A}.
Compare the total displayed curvature trajectories, including the
global terms, at $t_a$.  The discrepancy at $T_D$ has accuracy
$\min_{[j,D]}b_2>b(a)$; the difference of global bins is multiplied
by $O(s^a)$, and the residual $j$-slope by $O(s^j)$.
The resulting error is $s^{b(a)-o(1)}$, and the slope error is
$s^{\zeta-o(1)}$.  The raw derivative discrepancy fits the
displayed derivative accuracy by \eqref{eq:A}.  On the shared base
cell, the shear-difference Taylor formulas then give the normal
widths.  Hybrid fake drift fits by the strict slack and
\eqref{eq:Pl}.  These comparisons do not estimate the large
global terms separately at absolute time.  In hybrid, the forward
inequality \eqref{eq:Sp} transfers the already flattened $j$-row
scores to the final law; classically the comparisons are pathwise.
Thus the needed scores on the two systems agree, with power-small
strict-tolerance exceptions and persistence under restrictions of
vanishing exponent.  No earlier gate is modified after computing
later profiles.

This comparison concerns the canonical names and their angular
refinements.  It does not identify the new fine line-coefficient
side data at $j$ with analogous data at $D$.  Those finer data
are introduced on the $j$-law below, and are predicted from
earlier gates ending at $j$.  In particular the coarse bound
$K_j-K_D=O(s^{\zeta-o(1)})$ is not used as a depth-$30$
prediction of their normalized slopes.

Fix a typical $a<D$ for the boundary memory, differentiability,
and \eqref{eq:ng}.  Choose $a<j<D$, rational in hybrid, and then
$\delta=s^h$ with $h$ small compared with these fixed gaps.
For the latent side choose $A_0<a$, $V>v_0$ with the patch slack
of the coarse-side construction, and include
\[
 \mathcal C^{[j]}(A_0,V),\quad t_a,\quad
 R_a^{[j]}=[\beta'_j(T_j)]_{b(a)},\quad \bar k_j.
\]
Add depth-$30$ bins, in $\delta$ units, of the bounded quantities
\[
 (\beta'_j(t_a)-R_a^{[j]})s^{-b(a)},\qquad
 (K_j-\bar k_j)s^{-\zeta}.
\]
Use their representatives for $(\beta_F,\kappa_F)$ and choose the
representative root and phase on this refined side.  With
\[
 \theta=(T_j-t_a)s^{-a},\qquad
 r=(\beta'_j(T_j)-R_a^{[j]})s^{-b(a)},
\]
Equation~\eqref{eq:s-line} follows because $b(a)=a+\zeta$.

The conditional time-score bound follows from the causal time-list
count.  Indeed a row at any nearby prefix gate between $a$ and
$a+O(h)<j$, on a path to the tested event, predicts the side to
small lists.  Put $\Delta_K=b_2(j)-j-\zeta=(1-q)(D-j)>0$
and choose $100h<\Delta_K$, with $h$ also smaller by a fixed
large factor than all static-cell, derivative, and normal-drift
slacks.  On a connection from a gate $i\in[a,a+O(h)]$ to $j$,
Equation~\eqref{eq:A} gives
\[
 K_j-K_i=O(s^{\zeta+\Delta_K-o(1)}),\qquad
 \beta_j(t_a)-\beta_i(t_a)
   =O(s^{b(a)+\Delta_K-o(1)}).
\]
For the second bound use
$\min_{[i,j]}b_2\ge a+b_2(j)-j$ and multiply the slope
discrepancy by $O(s^a)$ when extrapolating to $t_a$.
Thus both normalized errors are finer than the new depth-$30$
bins.  The inequality
\[
 b_2(j)-j>\zeta
\]
controls both $K_j-K_i$ and the extrapolated curvature discrepancy
at $t_a$ more accurately than every new tangent bin, by
\eqref{eq:A}.  First subtract the guessed global bin.  The
residual $j$-slope has size $s^\zeta$, so the $T_j$-based root
curvature cut also has a small list.  Earlier static phase and
derivative data are predictable using
$l(A_0,V)<\min b_1$, exact base flow, and the hybrid normal slack.
The expanded lists at the prefix gate therefore exclude too-popular
time bins, with power-small exceptions at each fixed strict
tolerance.  In classical final-path tests the exact particle index
is retained throughout.

Let $S_i$ consist of depth-$12$ marks and quadratic fiber entries,
and put $O_i=\mathcal C^{[j]}(a+h,v_0)$.
The pair $(C,\mathcal C^{[j]}(a+40h,v_0))$ reads $S_i$ to
negligible lists.  Actual-root comparisons and the constant-cut
displayed-shear calculations give the output recoveries in
Lemma~\ref{act:line-test}.  The deterministic profile transfer
just proved transfers \eqref{eq:ng} to its stationarity hypothesis,
with the prefix lower bound from final memory; moreover
\[
 H_\delta(O_i\mid C)\le\partial_a\mathcal H(a,v_0)+o(1).
\]
These statements persist under a vanishing-exponent fraction
restriction of the $j$-row law.

We prove the remaining mismatch property while keeping $a,j,D,h$
fixed.  Suppose a $\delta^\chi$ interval for $\kappa_F$ is predicted
from $(C,S_i)$, where $0<\chi<1$ is fixed and small.  Choose
$0<\omega<\chi$ and group on $[j,D]$ by bins of $K_j$ and
$\beta_j(T_j)$ of widths
\[
 s^{\zeta+\omega h},\qquad s^{b(j)+\omega h}.
\]
Starting rows know these bins.  The end label knows the final
canonical names at the used times below $j$ at width $v_0$, hence
their $j$-versions, $C,S_i$, and $\bar k_j$, to small lists.
It therefore predicts the slope bins on success paths.  Given
one such slope, it guesses the value bin at $T_j$ from
$\beta_D(T_D)$: Equation~\eqref{eq:A} and time transport apply
because $\min_{[j,D]}b_2>b(j)$.  Include the ordinary time and
start-rounding choices.  After subtraction the residual end value
is bounded by $s^{b(j)+\omega h-o(1)}$, and the residual slope
times the terminal duration by $s^{b_2(D)-o(1)}$.
Since $b_2(D)-b(j)=D-j$, the curvature cost on $[j,D]$ is at most
$1-\omega h/(D-j)$, a fixed strict improvement.

Classically, failure of \eqref{eq:mis} yields a prediction with
nonvanishing final mass.  Use the Borel inner success and completion
marks of Lemma~\ref{cyl:borel-paths} at $D$, inside the current
end-state support;
Equation~\eqref{eq:Sp} and the count bounds supply the group lower
bound, using outer hulls of unrestricted geometric predecessors.
Exact particle and time starts fix the earlier gate
ownership, so these are actual path groups.  Group extraction
contradicts minimality of the curvature cost.

In hybrid the working law is $\mathbf P^j$, and a prediction-success
event there might have little terminal mass.  To address this,
discard before $j$-synthesis all rows whose depth-$\chi$ slope bin
has conditional mass at least $s^\tau$ given $(C,S_i)$, with
$\tau>0$ sufficiently small.  These maps use the single observation
on that row and their Borel versions under the full analyzed $j$-row
energy law
as in Lemma~\ref{cyl:borel-paths}.  For each conditioning value there are at most
$s^{-\tau}$ discarded bins.  We use a closure argument because
$\tau$ is not yet a vanishing exponent.  Hold
$a,j,D,h,\chi,\omega$ fixed, and put
$\gamma_0=\omega h/(D-j)>0$.  If arbitrarily small $\tau$
allowed the discarded synthesis to approach terminal saturation,
choose $\tau_n\downarrow0$ and genuine approximating experiments
whose terminal deficit and prior mesh, profile, and regularization
losses tend to zero faster than $\tau_n$.  For each $n$ group
on the same fixed interval $[j,D]$ by the preceding construction.
Each reachable end label has the expanded list of groups.
Finite-overlap inequalities have normalized loss
$O(\tau_n/(D-j))$, while other list and approximation losses
can be made $o(\tau_n)/(D-j)$.  Regularity is retained before
propagation, grouped energies sum to the $j$-energy budget,
and counts have the same overlap loss.  The quantitative loss
form of group extraction selects genuine experiments whose
extremality and count losses tend to zero in the $[j,D]$ units.
Their curvature costs are at most $1-\gamma_0+o(1)$, since
$h,j,D,\omega$ remain fixed.  They realize the original extremal
exponents in the closure with cost at most $1-\gamma_0/2$
eventually, contradicting minimality.  This applies the
vanishing-overlap clause along $\tau_n$; it does not apply a
zero-loss extraction statement at fixed $\tau$.
Consequently some sufficiently small fixed $\tau>0$ leaves a
strict positive terminal-exponent deficit on the discarded part.
The triangle inequality for the terminal
functional shows that the complementary synthesis preserves the
terminal exponent.  Its pre-synthesis row fraction is consequently
$s^{o(1)}$.

Restrict only the $j$-row probability law to this complement.
On retained rows the original conditional slope-bin mass is below
$s^\tau$.  Except on a vanishing restricted fraction, conditioning
normalization given $(C,S_i)$ multiplies mass by at most
$s^{-\tau/2}$: the excluded conditioning cells have original
retained mass at most $s^{\tau/2}$, whereas the total retained
mass is $s^{o(1)}$.  Neighboring-bin counting therefore makes
posterior collisions at the tested precision vanish.  The same
fixed representative maps, time bounds, and deterministic scores
are preserved by this restriction.  For finitely many separation
queries repeat the exclusion with previous errors sufficiently
small; diagonal limits give \eqref{eq:mis} for every fixed
threshold.  This is a restriction of the $j$-law for the entropy
test, not a claim about the final law of a modified stack.  The
strict curvature improvement is needed only for fixed $h$.

Lemma~\ref{act:line-test} now gives
$\partial_a\mathcal H(a,v_0)\ge2+3d$.  On this fixed-width
quadratic branch, $\partial_a W=1-\eta$ and
$p+d=1+3d-2\Gamma$ in \eqref{eq:canon}; hence
$\partial_a E\ge2\Gamma+\eta\ge2\Gamma$.
Finally, the inequalities used above hold for every $q<1$:
$b_2(j)-j-\zeta=(q-1)(j-D)>0$, and
$\min_{[j,D]}b_2>b_2(D)+j-D$ whether $q$ is positive or negative.
\end{proof}

\subsection{Finite motion when the angular trace may vanish}

Throughout the remaining argument write $k_0=1-d$ and
$p+d=1+3d-2\Gamma$, where $0<d\le1$.  Motion is in backwards
time $t=1-a$, at speed $k=dv/dt\ge0$.  Thus
\eqref{eq:surplus} makes $E$ decrease at rate at least $2\Gamma$.
On a fixed compact template domain the limiting profiles are
Lipschitz with a fixed constant, uniformly for small normal
enlargements $g>0$.  The two normal-bin list costs give
\begin{equation}
 \abs{E^g-E}\le Cg.
 \tag{enl}\label{eq:enl}
\end{equation}
Original and enlarged names may have deterministic profiles on
the same law: first request a countable sequence $g\downarrow0$
and dense depth lists; later finite requests are readout refinements.

\begin{lemma}[Finite off-gap motion]\label{act:finite-step}
Consider either an unclipped classical quadratic equality branch
with $b'=c>3$, $b_1'>0$, and uniform strong-prediction slack
$p_*(a)>b(a)+v$, or a classical linear branch with rate $c\ge1$.
The latter may be the right part of a clipped equality template,
or a pure linear template.  Fix a small $g>0$ and let $h>0$ be
sufficiently small compared with $g$ and the prediction slack.
Set $\delta=s^h$.  Let $\mathcal K$ be a sufficiently fine uniform
finite grid in $[c,c+1]$, including its endpoints, and set
\[
 k_R=k/2\quad\text{in the quadratic case},\qquad
 k_R=k\quad\text{in the linear case}.
\]
In particular $k_R\ge1$.  For a fixed $K_*$ sufficiently large
compared with $c+1$, define
\[
 C=\mathcal C^g(a,v),\quad O=\mathcal C^g(a+h,v),\quad
 C^k=\mathcal C^g(a,v+k_Rh),\qquad
 M_X=\lim H_s(X^{v+K_*h}\mid C).
\]
Provided the names stay on the specified branch, some $k\in\mathcal K$
satisfies
\begin{equation}
 E^g(a+h,v)-E^g(a,v+k_Rh)\ge\Gamma h-M_X.
 \tag{step}\label{eq:step}
\end{equation}
\end{lemma}

\begin{proof}
Use the actual $C$-shear at $t_a$.  Center the residual normal
velocity and starting position in their $C$-bins, and divide by
$s^{w-g},s^{a+w-g}$, respectively.  The bounded pair
$V=(V_n,V_d)$ is a function of $(y_0,C)$.  The output and root
know $\theta'=(t_{a+h}-t_a)/s^a$.  Given $C$, let $S_k$ contain
\begin{enumerate}[label=\textup{(\roman*)}]
\item $X^{v+K_*h}$;
\item $B_{t_a}$ to depth $a+v+k_Rh$ and
$Y_{t_a}-2[X]_{v+K_*h}B_{t_a}$ to depth $a+2v+2k_Rh$;
\item in the quadratic case,
$P-2R_aX+2k_aB_{t_a}$ to depth $l(a,v)+k_Rh$, in the root curvature cuts.
\end{enumerate}
These data refine with $k$, up to zero-cost lists.  Request the
finite deterministic profiles of $(S_k,V_z)$, $z,k\in\mathcal K$,
and their needed joints, without altering the earlier profiles.
Write $D(z\mid k)=\lim H_\delta(V_z\mid C,S_k)$.

Two geometric comparisons will be used:
\begin{equation}\label{act:step-geometry}
\begin{gathered}
 (O,C)\text{ predict }(V_n)_c
       \text{ and }(V_d+\theta'V_n)_{c+1},\\
 (C,S_k,V_k)\text{ predict }C^k,\qquad
 \lim H_s(S_k\mid O,C)\le M_X+h(3k_R-2).
\end{gathered}
\end{equation}
Here and below all prediction lists have vanishing exponent at
fixed $h$.  To verify the first line, the child basis minus the
root basis on the actual ray has derivative discrepancy $O(s^l)$
and curvature discrepancy $O(s^b)$ in $s^a$ duration units.
In curvature cases $b+v\ge l$.  Evaluating this difference at
the known root base-bin center instead costs only $O(s^w)$ in
normal velocity and $O(s^{a+w})$ in position, including the
vertical base tilt.  Since $(c+1)h\ll g$, these errors fit both
child accuracies in the enlarged normal units.  Evaluate all
known offsets at the exact bin times.  Thus the predictions are
single centers with the required radii, not merely unspecified
lists.  For $C^k$, the $S_k$ data supply its refined base and
derivative bins.  The same shear-difference comparison bounds the
normal error within its requested flat widths.

For the entropy bound in \eqref{act:step-geometry}, first give
$X^{v+K_*h}$.  Backflow from $(O,C)$ reads both the horizontal and
tilted vertical positions at normalized depth at least $1$.
In the tilted coordinate the velocity is
$X^2-2[X]_{v+K_*h}X$, whose error about the refined center is
quadratic.  The two extra position costs are therefore
$(k_R-1)h$ and $(2k_R-1)h$.  In the quadratic case the output
displays derivative refinement $ch$, which suffices because
$k_R\le c$.  Changing back to the root curvature basis costs
$O(s^b)$ in earlier-time units; multiplying by the unresolved
base errors in the guessed position bins costs at most
$O(s^{l+k_Rh})$.  No further derivative entropy is required.

Conditionally on $(y_0,C,S_k)$, time bins at $\delta$-depth
$z\in[0,1]$ have typical score at least $dz$.
Indeed $(C,S_k)$ is predicted to small lists from $(y_0,t_a)$ on
all paths outside the profile exceptions.  Root curvature
predictions follow from \eqref{eq:A}, \eqref{eq:F}, and cut
ownership.  The quadratic derivative query is strictly below
$p_*(a)$; in a clipped linear branch it uses only $l=b_1(a)$,
and in pure linear form it uses the center slope comparison.
Exact classical flow predicts the phase entries.  Expanding any
popular time list by these predicted lists and applying the causal
time count through $a+zh$ gives power-small failures at each fixed
strict tolerance.  Quantizing the child bin-start time changes
only bounded factors.  Likewise, deterministic profiles and the
neighboring-bin likelihood bound give conditional $V$-ball mass
at most $\delta^{D(z\mid k)-o(1)}$ about typical samples, for
$z\in\mathcal K$.  These are bounds for the unrestricted
conditional mass of those balls.  The trimmed time law has
unnormalized interval bounds of exponent $d$.

We claim
\begin{equation}
 H_\delta(O\mid C,S_k)
 \ge d+\min_{z\in\mathcal K}
       \{D(z\mid k)+d(c+1-z)\}
       -O(\operatorname{mesh}\mathcal K).
 \tag{sh}\label{eq:sh}
\end{equation}
The first term is supplied by $H_\delta(O\mid y_0,C,S_k)$.
For the information with $y_0$, draw an independent reference
$y'_0$ given $(C,S_k)$, and let $V'$ be its normal pair.  If it
is admissible for $(O,C)$, the first comparison in
\eqref{act:step-geometry} forces agreement in $V_n$ to
$O(\delta^c)$ and in $V_d+\theta'V_n$ to
$O(\delta^{c+1})$.  Hence $\abs{V'-V}\lesssim\delta^c$.
If its distance lies between consecutive grid widths
$\delta^{z+\mu}$ and $O(\delta^z)$, where
$\mu=\operatorname{mesh}\mathcal K$, the velocity difference
is comparable to the full difference unless the latter is already
at final width.  Consequently the time projection confines
$\theta'$ to an interval of length
$O(\delta^{c+1-z-\mu})$.  At the final width no time gain is
charged.

Average reference admissibility over actual samples satisfying
both typicality bounds.  Given the actual $(y_0,C,S_k)$, use the
conditional ball bound for $V'$ and, for each fixed $V'$, the
trimmed time interval bound.  Summing the finitely many annuli
bounds this average by the power in \eqref{eq:sh} minus its first
$d$, with strict tolerance and the $O(\mu)$ loss.  Markov's
inequality makes the reference admissibility probability at most
that power with slightly more tolerance for most actual
$(O,C,S_k)$.  The actual posterior given $(O,C,S_k)$ is supported
on the admissible set.  The relative entropy bound for a measure
supported on a reference set of mass $q$, namely $\log(1/q)$,
therefore gives the information term in \eqref{eq:sh}.  All
trimming was on the actual sample in this averaging argument;
no exceptional-set transfer to the independent product law is used.

For $z_2>z_1$, the increment
$D(z_2\mid k)-D(z_1\mid k)$ is nonincreasing in $k$, because
$S_k$ refines.  Thus the largest minimizer $z(k)$ in
\eqref{eq:sh} is nondecreasing.  A nondecreasing self-map of a
finite ordered set has a fixed point: iterate from its least
element.  At such a point $z(k)=k$, use
\eqref{act:step-geometry} and refinement of $C$ by $O$ to obtain
\[
 \lim\{H_\delta(O)-H_\delta(C^k)\}
 \ge d+d(c+1-k)+2-3k_R-M_X/h-O(\mu).
\]
The weight exponent difference, later minus moved, divided by $h$,
is $-(3+\eta)k_R+(1-\eta)(c-k)$.  Subtracting this and $p+d$
as in \eqref{eq:canon} gives
\[
 \frac{E^g(a+h,v)-E^g(a,v+k_Rh)}h
 \ge(1-d)(1+k-c)+2\Gamma
      +\eta(k_R+c-k)-M_X/h-O(\mu).
\]
Both extra displayed terms are nonnegative: $k\in[c,c+1]$,
$d\le1$, and $k_R+c-k\ge0$ in the stated branches.
Choose $\mu$ small enough to retain $\Gamma$.
The finite refinement proof can be applied at any macroscopic
path node to the same profiles $E^g,M_X$; it does not select
microscopic configurations adaptively.  This proves \eqref{eq:step}.
\end{proof}
\subsection{The off-gap path contradiction}

The finite step loses only the angular information $M_X$. We compare that
loss with angular scores at later backward times, so that its accumulated
contribution vanishes across widths where the trace is zero. This lets us
join finite steps to the differential motions from the interior estimates.

Write $E_t^g(v)=E^g(1-t,v)$, $E_t(v)=E_t^0(v)$, and
$m_t(v)=m(1-t,v)$.  Denote the cost in \eqref{eq:step} with
earlier root $(1-t-h,v)$ by $M_X(t,v;h,g)$.  Thus a permitted
speed $k$ satisfies
\[
 E_t^g(v)-E_{t+h}^g(v+kh)\ge\Gamma h-M_X(t,v;h,g).
\]
On a fixed compact canonical domain there is a constant $A_0$ such
that, if $t'>t+h+A_0g$ and $0\le v-h\le y\le v$, then
\begin{equation}
 M_X(t,v;h,g)
 \le U\bigl(y+(K_*+1)h;1-t',y\bigr).
 \tag{shift}\label{eq:shift}
\end{equation}
Indeed the enlarged name at $(1-t-h,v)$ predicts the ordinary
name at $(1-t',y)$.  Nesting handles time, curvature, derivative,
and base data.  The positive forward rate of the normal-width exponent
absorbs $g$ in this time shift; backflow in the fine shear costs
its enlarged velocity width and that width times the longer
duration, and the shear-difference estimate on the base cell
fits the ordinary coarse widths.  This remains valid across a
clipping switch.  Finally $y+(K_*+1)h\ge v+K_*h$, so conditioning
and refinement prove \eqref{eq:shift}.  For each fixed rational
$t'$ the angular profile properties give
\begin{equation}\label{act:shift-limit}
 F_h(t',y):=h^{-1}U(y+(K_*+1)h;1-t',y)
 \longrightarrow(K_*+1)m_{t'}(y)
\end{equation}
for almost every $y$, with $0\le F_h\le K_*+1$.

\begin{lemma}[Off-gap path contradiction]\label{act:off-gap-path}
Suppose that on an open time interval and in a sufficiently narrow
right neighborhood of $v=Jt$ the following hold.
The profiles $E,E^g$ are nonnegative and uniformly Lipschitz on the
compact domains used, satisfy \eqref{eq:enl}, and $E_t(Jt)=0$.
Equations \eqref{eq:step} and \eqref{eq:shift} apply, with bounded
step speeds whose minimum is at least $1$ and strictly greater
than $J$; step admissibility is uniform on compact subsets for
$h/g$ small enough.  The angular rate has its nondecreasing-in-$t$
version.  There is a speed $k_+\ge1$, $k_+>J$, for which
\eqref{eq:surplus} holds almost everywhere where $m_t(v)>0$.
Alternatively, that motion need only hold for $0<m_t(v)\le k_0$
provided the vertical estimate $\partial_tE_t(v)\le-2\Gamma$
holds almost everywhere where $m_t(v)>k_0$, throughout the
constant-width continuations used, including after crossing the
seam.  The normal-width exponent has positive forward rate on
all branches needed for \eqref{eq:shift}.
These conditions are incompatible with $\Gamma>0$.
\end{lemma}

\begin{proof}
Fix compact domains and all Lipschitz and speed bounds before
choosing density scales.  Let $C_0$ denote constants depending
only on these bounds and $K_*$.  Outside a null set of widths set
\[
 \sigma(v)=\inf\{t\in\mathbb Q:m_t(v)>0\},
\]
with the terminal endpoint convention if the set is empty.
Take rational queries in a containing interval and truncate this
threshold to its closure if necessary.  Monotonicity implies
positivity at almost every relevant point above $\sigma$, whereas
$m_{t'}(v)=0$ at each rational $t'<\sigma(v)$ outside its null
set of exceptional widths.

We first justify integration on a prescribed positive segment.
Fix a slab $[s,T]$ and a measurable width set $P$ where $m_t(v)>0$
for almost every $t$ in that slab.  Consider a proposed segment
of speed $k_+$, duration $\ell$, and width interval $W$.
At its initial point one has either
\begin{align}
 E_s(v)&\ge2\Gamma\ell+E_{s+\ell}(v+k_+\ell)
                  -C_0|W\setminus P|,
                  \label{act:positive-segment}\\
 \text{or}\qquad
 E_s(v)&\ge2\Gamma(T-s)-C_0|W\setminus P|.
                  \label{act:terminal-freeze}
\end{align}
To see this without assuming the proposed line is generic, perturb
its initial width.  For almost every perturbation, Fubini's theorem
and the Lipschitz chain rule give the moving estimate almost
everywhere and valid vertical estimates at almost every encountered
width.  Time spent outside $P$ is bounded by
$|W\setminus P|/k_+$, with an error tending to zero under the
perturbation.  If freezing is allowed and the high-rate set
$m>k_0$ on that line has positive time measure, choose a regular
point arbitrarily close to its first essential occurrence.
The measure of earlier high-rate times can be made arbitrarily
small.  Integrate the moving estimate to that point, using the
Lipschitz bound on those earlier high-rate times, and then freeze
to $T$.  Monotonicity retains $m>k_0$ on the vertical continuation,
and endpoint nonnegativity proves \eqref{act:terminal-freeze}.
If there is no positive-measure high-rate set, integrate to the
moving endpoint to get \eqref{act:positive-segment}.
Let the earlier-high-time tolerance and then the perturbation
tend to zero.  A subsequence either always freezes or always
reaches the proposed endpoint, and Lipschitz continuity gives the
corresponding alternative at the original starting point.

Suppose $J>0$ and put $T(v)=v/J$ on an interior seam-width interval.
It is covered by
\[
 Z_0=\{\sigma(v)\ge T(v)\},\qquad
 \{\sigma(v)\le T(v)-1/n\},\quad n\ge1.
\]
One of these sets has a density point in that interval.
If $v_*=Jt_*$ is a density point of
$P=\{\sigma(v)\le T(v)-\Delta\}$, $\Delta>0$, fix a small
$\eps>0$ and start at $(t_*,v_*+\eps r)$.  Propose a positive
segment of duration $r$.  For $r$ small, every visited width in
$P$ has threshold below $t_*$, while density gives
$|W\setminus P|=o(r)$.  The segment stays to the right of the
seam.  Either \eqref{act:positive-segment} or
\eqref{act:terminal-freeze}, with $T=t_*+r$, implies
\[
 E_{t_*}(v_*+\eps r)\ge2\Gamma r-o(r).
\]
Seam equality bounds the left side by $C_0\eps r$, a contradiction
after first choosing $\eps$ small and then $r$ small.

Now let $v_*=Jt_*$ be a density point of $Z_0$.  Start at the same
width and make steps of duration $h$ until less than $h$ remains
in $[t_*,t_*+r]$.  At a step starting at $(t_i,v_i)$, the speed
bound gives $v_i-Jt_i\ge\eps r$.  For $h$ small compared with
$\eps r$, all $y\in I_i=[v_i-h,v_i]$ satisfy
\[
 T(y)-(t_i+h)\ge b_0:=\frac{\eps r}{2J}.
\]
These intervals have disjoint interiors, since every step
increases width by at least $h$.
Choose a fixed finite rational grid meeting every interval
$(t+b_0/4,t+b_0/2)$ for $t\in[t_*,t_*+r]$.  Choose $g$ with
$A_0g<b_0/4$.  For each step take a grid point $t'_i$ in this
interval for $t=t_i+h$.  It is admissible in \eqref{eq:shift}
for the entire $I_i$.  On $I_i\cap Z_0$ it satisfies
$t'_i<T(y)\le\sigma(y)$, and hence $m_{t'_i}(y)=0$.
Integrating \eqref{eq:shift} divided by $h$ over $I_i$ gives
\[
 M_X(t_i,v_i;h,g)\le\int_{I_i}F_h(t'_i,y)\,dy.
\]
On the good portions this integrand is dominated by
\[
 \sum_{t'\text{ in the fixed grid}}
       \mathbf1_{\{t'<\sigma(y)\}}F_h(t',y),
\]
which tends to zero almost everywhere and is bounded.  Notice
that only grid times below the threshold are included.
All visited widths lie in a window $W_r$ of length $O(r)$ based
at $v_*$.  Disjointness and bounded convergence give
\[
 \sum_i M_X(t_i,v_i;h,g)
 \le C_0|W_r\setminus Z_0|+o_{h\to0}(1).
\]
First choose $r$ using density so that the first term is small
compared with $\Gamma r$, then $g\ll\eps r$, and only then
let $h\to0$.  Telescoping \eqref{eq:step}, using
\eqref{eq:enl} and endpoint nonnegativity, gives
\[
 \Gamma(r-h)\le C_0\eps r+C_0g
                  +C_0|W_r\setminus Z_0|+o_{h\to0}(1),
\]
which is impossible with these choices.

Suppose finally $J=0$.  Fix $\eps>0$ and a compact positive-length
interval $I_0$ of starting times.  Let $L_0$ exceed the speed bound
by a fixed margin.  Fubini's theorem gives
\[
 \int_{I_0}\bigl|\{0<v<L_0r:
          \sigma(v)\in[t-r,t+2r]\}\bigr|\,dt
 \le3L_0r^2.
\]
For $r$ sufficiently small, there is therefore $t_*\in I_0$ with
transition-set length at most $\eps r$.  In $(0,L_0r)$ put
\[
 P=\{\sigma(v)<t_*-r\},\qquad
 Z_0=\{\sigma(v)>t_*+2r\}.
\]
The remaining set has length at most $\eps r$.  Widths in $P$
have positive trace throughout the slab $[t_*,t_*+r]$.
Choose a rational $t'\in(t_*+5r/4,t_*+7r/4)$; the trace vanishes
at that time for almost every width in $Z_0$.
Approximate $P$ to symmetric-difference error $\eps r$ by a
finite union $P'$ of intervals, with $N$ endpoints.  Now choose
$g$ with $A_0g<r/8$ and $C_0(N+1)g\le\eps r$, and subsequently
let $h\to0$ with $h\ll g$ and $h<\eps r/2$.

Start at $(t_*,\eps r)$.  Outside $P'$ make an off-gap step;
inside $P'$ apply the positive-segment comparison up to its next
right endpoint or the final time.  Stop immediately if the
terminal alternative \eqref{act:terminal-freeze} occurs.
Consistent half-open endpoint conventions prevent zero-length
transitions.  Width increases before any terminal freeze, so
there are at most $N+1$ positive segments.  Changes between
$E$ and $E^g$ cost at most $C_0(N+1)g$.
Their positive-segment width intervals are disjoint and lie in
$P'$, so the total bad-width charge is at most
$C_0|P'\setminus P|\le C_0\eps r$.

For off-gap steps the intervals $I_i=[v_i-h,v_i]$ again have
disjoint interiors.  Since their right endpoints are outside
$P'$, their intersections with $P'$ lie within distance $h$ of
its endpoints.  Consequently
\[
 \left|\left(\bigcup_iI_i\right)\cap P\right|
 \le\eps r+2Nh.
\]
The one fixed $t'$ is admissible in \eqref{eq:shift} for every
step.  On $Z_0$, Equation~\eqref{act:shift-limit} and bounded
convergence give zero limiting integral; on the transition set
and the parts of $P$ just estimated use $F_h\le K_*+1$.
Hence
\[
 \sum_iM_X(t_i,v_i;h,g)
 \le C_0\eps r+C_0Nh+o_{h\to0}(1).
\]
Each off-gap step pays $\Gamma$ times its duration and each
positive segment pays $2\Gamma$, subject to the listed errors.
A terminal freeze pays to the final time; otherwise less than
$h$ remains unused.  Initial seam equality, initial displacement,
and final nonnegativity yield
\[
 \Gamma(r-h)\le C_0\eps r+C_0(N+1)g+C_0Nh+o_{h\to0}(1).
\]
Let $h\to0$ with earlier choices fixed.  The choice of $g$ gives
a contradiction for $\eps$ sufficiently small, completing the proof.
\end{proof}

\subsection{Exhaustion of the affine configurations}

\begin{proof}[Proof of Theorem~\ref{thm:cylinder}]
First suppose a classical extremal configuration has
$\Gamma>\eta/2$.  After it is excluded, suppose a hybrid
configuration has $\Gamma>\eta/2$.  It then has the strict gap
$\Gamma>\max(\Gamma_{\mathrm{cl}},\eta/2)$ required for
\eqref{eq:Pl}.  The affine reduction \eqref{eq:F} and the
zero-cost arguments leave precisely the positive-cost cases
treated below.  All uses of \eqref{eq:surplus} are on valid
branches with the indicated crossing or spreading hypotheses.

We describe once the integration conventions.  Almost-everywhere
estimates on affine rays are integrated on generic arbitrarily
nearby rays and passed to the desired one by Lipschitz continuity.
A speed with surplus $2\Gamma$ may be decreased slightly while
retaining surplus at least $\Gamma$.  Rays from the terminal
corner can start at a small positive backward time, with initial
$E$ tending to zero.  If a high-rate linear portion is encountered,
the freezing alternative in the proof of
Lemma~\ref{act:off-gap-path} applies.

We also require a short-slab observation.  Suppose $b_1'>0$ and
choose a typical $j$ with $E(j,u(j))=0$.  Work in a short slab
ending at $j$, omitting a small fixed fraction next to its final
edge.  For $v>u(j)$ sufficiently close to $u(j)$, the pair
$(\lambda_j,X^v)$ recovers $\mathcal C(a,v)$ to small lists
throughout the shortened slab.  In the unclipped quadratic case,
choose it so that $b(a)+v<b_1(j)$.  If $R_j$ is the angular
profile given $\lambda_j$, conditioning gives
\begin{equation}\label{act:slab-trace}
 m(a,v)\ge R_j'(v)
 \quad\text{for almost every such }(a,v).
\end{equation}
An integrated lower slope for $R_j$ through $u(j)$ supplies
positive-measure sets of arbitrarily close widths with derivative
at least that slope minus any prescribed strict error.  Otherwise
integration would contradict the lower slope.  A bound above
capacity is already impossible.  Choose these widths generically
for all almost-everywhere assertions.  If the supplied derivative
exceeds a freezing threshold, \eqref{eq:surplus} holds throughout
the slab, except for any stipulated small time-density error.
Its positive integral contradicts nonnegativity: the terminal
width perturbation and the omitted edge cost arbitrarily little
by Lipschitz continuity and $E(j,u(j))=0$.
When a threshold depends on the mixed trace $\nu$, choose $j$
to be a Lebesgue point of that trace.  The strict margin at $j$
then survives off a time set of arbitrarily small relative length.
All extra $j$-profile queries are readout refinements after $j$
has been selected.

Finally, if $p_*>b_1$ on a positive-measure set of interior times,
monotonicity of $p_*$ and continuity of $b_1$ give a smaller open
interval with uniform strict prediction slack.  Otherwise
$p_*=b_1$ at almost every time under consideration.  We call
these the strong and weak alternatives, respectively.

\paragraph{Equality curvature: $q=c_2>0$.}
Put $c=c_2$ and $C_1=c-J=b_1'$.  If $J>0$, consider left
quadratic points $0<v<u(a)$ with
$l<\min_{[a,1]}b_1$ and, in hybrid, $w<0$.
Write $v=u(j)$ for a later $j<1$.  Left memory gives
\begin{equation}\label{act:left-memory}
 m\ge\frac{k_0+2\Gamma+\eta(L_{\mathrm{raw}}-1)}
               {L_{\mathrm{raw}}}>0,
 \qquad
 L_{\mathrm{raw}}=\max\{1,(1+c)/2,1+c-J\}.
\end{equation}
The coarse-side construction and avoidance supply radial spreading,
including the extra left-slack version when $c=1$.
If $0<c\le1$, then $L_{\mathrm{raw}}\le1+c$, so
\eqref{act:left-memory} exceeds the quadratic freezing threshold
$k_0/(1+c)$.  On every compact interior time slab all sufficiently
small $v>0$ meet the hypotheses.  Integrate at those widths,
let $v\downarrow0$, and then let the trimmed endpoint tend to
$1$.  The equality $E(1,0)=0$ contradicts the resulting positive
drop.  The same argument applies for $c>1$ when
$L_{\mathrm{raw}}\le2$, using the threshold $k_0/2$.

If $c>1$ and $J\ge(c+1)/4$, use a ray from the terminal corner
at a speed just below $(c+1)/4$.  It stays in the left region,
has $2k<c$, and satisfies the derivative prediction condition;
if $b_1$ decreases forward, $l<0$ suffices.  The positive
quadratic motion in \eqref{eq:surplus} is contradictory.
This covers hybrid equality configurations because
\eqref{eq:Pl} gives $J=(c+1/2)/2$ there.

The remaining configurations are classical and have $b_1'>0$:
either $J=0$, or $c>1$ and
$0<J<(c+1)/4$, with $C_1>1$.  Indeed if $C_1\le1$ in the
latter range, then $L_{\mathrm{raw}}\le2$, a case already excluded.
In a strong interval use the unclipped quadratic template near
the seam with exact side and prediction slack.  Avoidance supplies
radial spreading; at $c=1,J=0$ this is supplied by
Lemma~\ref{act:stationary-boundary}.  If $c\le3$, the strengthened
memory bound \eqref{eq:mr} uses $L=\max\{1,(1+c)/2\}$.
For $c<1$ it exceeds $k_0/(1+c)$, and for $1\le c\le3$ it
exceeds $k_0/2$.  Equation~\eqref{act:slab-trace} therefore
gives the short-slab freezing contradiction.  Shrink the slab
to retain $l<p_*$ throughout.  If $c>3$, apply
Lemma~\ref{act:off-gap-path}: its positive speed is
$k_+=(c+1)/4>J$, and the finite-step speeds have minimum
$c/2>J$.  They are all at least $1$, the normal-width exponent increases
forward, and no high-rate switch is needed.

In the weak alternative use the clipped template.  Its right
branch $v>u(a)$ is linear of rate $C_1$, with exact construction
and crossing at almost every weak time.  If $c<1$, necessarily
$J=0$ here and $C_1=c$.  At a typical mixed-trace time $j$, use
both memory inequalities \eqref{eq:mr}.  Together with
\eqref{act:slab-trace}, their strict margin gives positive $m$
above $k_0-c\nu(a)$ outside an arbitrarily small relative set
of times.  The low-rate linear freezing bound contradicts the
short-slab observation.  In all other weak cases $C_1\ge1$.
For $0<m\le k_0$ use $k_+=(1+C_1)/2$, and finite-step speeds
have minimum $C_1$; both are strictly greater than $J$ in the
remaining range.  For $m>k_0$, freeze on the right by the linear
bound.  A continued freeze crossing left uses the quadratic
freezing bound, the same monotone clipped-template trace,
coarse side, and line avoidance.  The forward width rates are
$C_1$ on the right and $c$ on the left, both positive.
Lemma~\ref{act:off-gap-path} applies and gives a contradiction.

\paragraph{Pure linear cost: $c_2=0<c_1$.}
Write $c=c_1$.  The linear construction has crossing.
If classical and $c<1$, use both inequalities \eqref{eq:mr}
at a typical mixed-trace time and the short-slab argument just
given; there is now no side-of-seam restriction.
In hybrid, \eqref{eq:Pl} gives $J=c+1/2$.  Use a ray from the
corner at a speed slightly below $c$ for $c<1$, and slightly
below $(1+c)/2$ for $c\ge1$.  It lies in $v<u(a)$, $w<0$,
and left memory gives
\[
 m\ge\frac{k_0+2\Gamma+\eta c}{1+c}>0.
\]
For $c<1$ this exceeds the threshold $k_0/(1+c)$ for linear
motion at speed $c$; for $c\ge1$ use the stated positive
linear motion while $m\le k_0$.  If $m>k_0$ on a positive
time set, switch at a regular point to freezing.  Its hypotheses
remain valid backwards.  These motions contradict nonnegativity
and terminal equality.  The same argument, with exact side,
works classically for $c\ge1$ and $J\ge(1+c)/2$.
For the remaining classical case $c\ge1$, $J<(1+c)/2$, use
Lemma~\ref{act:off-gap-path} on the pure linear template:
$k_+=(1+c)/2>J$, finite-step speeds are at least $c>J$,
and any high-rate freeze stays in a linear branch.

\paragraph{Strict curvature: $q<1=c_2$.}
Put $C_1=q-J<1$ and use the artificial template ending at $D$.
At $v_0=u(D)$ memory supplies the lower slope
\[
 \frac{k_0+2\Gamma+\eta(L_{\mathrm{raw}}-1)}{L_{\mathrm{raw}}},
 \qquad L_{\mathrm{raw}}=1+\max(0,C_1).
\]
If $C_1\le0$, then $L_{\mathrm{raw}}=1$, so this is positive
and at least $1-2d$.  This explicitly includes every negative
$q$ allowed by the affine reduction.  If $C_1>0$, the same
conclusion follows whenever
$d\ge C_1/(1+2C_1)$, since
$(1-d)/(1+C_1)\ge1-2d$ is exactly that inequality.
Take $D=1$, integrate Proposition~\ref{act:strict-boundary},
and contradict $E(D,v_0)=0$.

It remains that $C_1>0$ and $d<C_1/(1+2C_1)$.
If $C_1\le1/2$, then $d<1/4$.  From the terminal corner
$D=1$ take a quadratic ray of speed slightly below $1/2$.
It satisfies $l<\min_{[a,D]}b_1$ and $w<0$ strictly in the
interior: along a speed $k<1/2$ ray these are
$-(1-k)(1-a)<-C_1(1-a)$ and
$-(1-2k)(1-a)<0$.  The coarse side is valid.
The unit-cost quadratic motion bound without spreading, which
also permits $m=0$, gives the contradiction.
All hybrid strict configurations have been covered:
\eqref{eq:Pl} gives $J=(q+1/2)/2$, so
$C_1=q/2-1/4<1/4$.

Finally suppose $1/2<C_1<1$.  The configuration is classical,
and $2d<C_1$.  If the raw final derivative has a strong interval,
choose a typical $D$ in it, retaining prediction slack immediately
to its left.  The strengthened memory estimate \eqref{eq:mr}
has $L=1$ for this boundary comparison via the encountered
$\lambda_D$ on final paths.  Proposition~\ref{act:strict-boundary}
again contradicts $E(D,u(D))=0$.
Otherwise take the clipped artificial template with $D=1$ and
move from the corner at speed $1$, using generic small
perturbations.  This is strictly right of the clipping switch
$v=b_1(a)-b(a)=(1-C_1)(1-a)$.  The classical linear branch has
crossing at almost every weak time.  Its arbitrary-$m$ motion
bound for $2d<C_1<1$, including $m=0$, gives the contradiction
without requiring a positive angular trace.

Every affine positive-cost configuration has now been excluded.
Together with the zero-cost arguments this proves the classical
bound, and subsequently the hybrid bound with its required
strict separation, as claimed.
\end{proof}
